% GENERATED FILE. Edit canonical lessons, then run npm run generate:books.
% canonical-source-hash: fnv1a64:768d9080228bfabf
% canonical-lessons: GU-C254-sangit, GU-C254-git, GU-C254-vagadvu, GU-C254-vasli, GU-C254-tablu, GU-R254-review-249-254

\chapter{Music and song}
\label{ch:gu-music-and-song}
\hlchaptermodality{254}

\begin{chapteropening}
\textbf{By the end of this chapter:} \emph{I can say what music I played and which songs I sang.}

Complete the last lesson of chapter 254: I can say what music I played and which songs I sang.
\end{chapteropening}

\section[\texorpdfstring{\gu{સંગીત} (sangīt)}{સંગીત (sangīt)}]{\gu{સંગીત} (sangīt) — music}

\label{lesson:GU-C254-sangit}



\begin{quote}
Before the new one: say the Gujarati for to get bored, then the Gujarati for to try out.
\end{quote}

\subsection*{You'll want to know: \gu{સંગીત}}
\textbf{\gu{સંગીત}} — \emph{sangīt} — \textquotedblleft{}music\textquotedblright{}.

\begin{quote}
\textbf{\gu{ગઈકાલે} \gu{અમે} \gu{સંગીત} \gu{સાંભળ્યું}.} — \emph{gaīkāle ame sangīt sā̃bhaḷyũ} — \textquotedblleft{}Yesterday we listened to music.\textquotedblright{}
\end{quote}

A neuter noun.

\begin{cousinweb}
\textbf{\gu{સંગીત}} (\emph{sangīt}) is Sanskrit \emph{saṁgīta}, \textquotedblleft{}chorus, concert\textquotedblright{}, taken over whole as a tatsama (CDIAL 12849); EtymDB also lists it as borrowed from Sanskrit.
\end{cousinweb}

\begin{grammarlens}[title={What you've built}]
One more word for everyday life at A2.
\end{grammarlens}

\subsection*{Your turn}
\begin{itemize}
\raggedright
  \item \emph{Say it:} \emph{sangīt}
  \item \emph{Say it:} \emph{sangīt}, once more
  \item \emph{Say it:} \emph{gaīkāle ame sangīt sā̃bhaḷyũ} — Yesterday we listened to music.
  \item \emph{Recall:} say the Gujarati for to get bored, then the Gujarati for to try out, then say \emph{sangīt} again
\end{itemize}

\begin{tcolorbox}[breakable,colback=teal!4,colframe=teal!35!black,title={Before you move on}]
What does \gu{સંગીત} mean? (\textquotedblleft{}Music\textquotedblright{}.) Say it once more.
\end{tcolorbox}
\section[\texorpdfstring{\gu{ગીત} (gīt)}{ગીત (gīt)}]{\gu{ગીત} (gīt) — song}

\label{lesson:GU-C254-git}



\begin{quote}
Before the new one: say the Gujarati for to try out, then the Gujarati for music.
\end{quote}

\subsection*{You'll want to know: \gu{ગીત}}
\textbf{\gu{ગીત}} — \emph{gīt} — \textquotedblleft{}song\textquotedblright{}.

\begin{quote}
\textbf{\gu{તેણે} \gu{એક} \gu{સુંદર} \gu{ગીત} \gu{ગાયું}.} — \emph{teṇe ek sundar gīt gāyũ} — \textquotedblleft{}She sang a beautiful song.\textquotedblright{}
\end{quote}

\begin{cousinweb}
\textbf{\gu{ગીત}} (\emph{gīt}) is Sanskrit \emph{gītá}, \textquotedblleft{}singing, song\textquotedblright{}, taken over as a tatsama (CDIAL 4167). The taught \emph{gāvũ}, \textquotedblleft{}to sing\textquotedblright{}, comes from the same root.
\end{cousinweb}

\begin{grammarlens}[title={What you've built}]
One more word for everyday life at A2.
\end{grammarlens}

\subsection*{Your turn}
\begin{itemize}
\raggedright
  \item \emph{Say it:} \emph{gīt}
  \item \emph{Say it:} \emph{gīt}, once more
  \item \emph{Say it:} \emph{teṇe ek sundar gīt gāyũ} — She sang a beautiful song.
  \item \emph{Recall:} say the Gujarati for to try out, then the Gujarati for music, then say \emph{gīt} again
\end{itemize}

\begin{tcolorbox}[breakable,colback=teal!4,colframe=teal!35!black,title={Before you move on}]
What does \gu{ગીત} mean? (\textquotedblleft{}Song\textquotedblright{}.) Say it once more.
\end{tcolorbox}
\section[\texorpdfstring{\gu{વગાડવું} (vagāḍvũ)}{વગાડવું (vagāḍvũ)}]{\gu{વગાડવું} (vagāḍvũ) — to play an instrument, to ring}

\label{lesson:GU-C254-vagadvu}



\begin{quote}
Before the new one: say the Gujarati for music, then the Gujarati for song.
\end{quote}

\subsection*{You'll want to know: \gu{વગાડવું}}
\textbf{\gu{વગાડવું}} — \emph{vagāḍvũ} — \textquotedblleft{}to play an instrument, to ring\textquotedblright{}.

\begin{quote}
\textbf{\gu{તેણે} \gu{ઢોલ} \gu{વગાડ્યો}.} — \emph{teṇe ḍhol vagāḍyo} — \textquotedblleft{}He played the drum.\textquotedblright{}
\end{quote}

\begin{cousinweb}
\textbf{\gu{વગાડવું}} (\emph{vagāḍvũ}), to make something sound, is the causative of the inherited \emph{vāgvũ}, \textquotedblleft{}to sound\textquotedblright{}. CDIAL lists \emph{vāgvũ} under Sanskrit \emph{vādyatē}, \textquotedblleft{}is made to sound\textquotedblright{} (CDIAL 11513).
\end{cousinweb}

\begin{grammarlens}[title={What you've built}]
One more word for everyday life at A2.
\end{grammarlens}

\subsection*{Your turn}
\begin{itemize}
\raggedright
  \item \emph{Say it:} \emph{vagāḍvũ}
  \item \emph{Say it:} \emph{vagāḍvũ}, once more
  \item \emph{Say it:} \emph{teṇe ḍhol vagāḍyo} — He played the drum.
  \item \emph{Recall:} say the Gujarati for music, then the Gujarati for song, then say \emph{vagāḍvũ} again
\end{itemize}

\begin{tcolorbox}[breakable,colback=teal!4,colframe=teal!35!black,title={Before you move on}]
What does \gu{વગાડવું} mean? (\textquotedblleft{}To play an instrument, to ring\textquotedblright{}.) Say it once more.
\end{tcolorbox}
\section[\texorpdfstring{\gu{વાંસળી} (vā̃sḷī)}{વાંસળી (vā̃sḷī)}]{\gu{વાંસળી} (vā̃sḷī) — flute}

\label{lesson:GU-C254-vasli}



\begin{quote}
Before the new one: say the Gujarati for song, then the Gujarati for to play an instrument.
\end{quote}

\subsection*{You'll want to know: \gu{વાંસળી}}
\textbf{\gu{વાંસળી}} — \emph{vā̃sḷī} — \textquotedblleft{}flute\textquotedblright{}.

\begin{quote}
\textbf{\gu{મેં} \gu{કાલે} \gu{વાંસળી} \gu{વગાડી}.} — \emph{mẽ kāle vā̃sḷī vagāḍī} — \textquotedblleft{}I played the flute yesterday.\textquotedblright{}
\end{quote}

\begin{cousinweb}
\textbf{\gu{વાંસળી}} (\emph{vā̃sḷī}) is a tadbhava. It comes from Sanskrit \emph{vaṁśī}, \textquotedblleft{}flute, pipe\textquotedblright{}, built on \emph{vaṁśa}, \textquotedblleft{}bamboo\textquotedblright{}; CDIAL lists Gujarati \emph{vā̃sḷī}, \textquotedblleft{}pipe, flute\textquotedblright{} (CDIAL 11180).
\end{cousinweb}

\begin{grammarlens}[title={What you've built}]
One more word for everyday life at A2.
\end{grammarlens}

\subsection*{Your turn}
\begin{itemize}
\raggedright
  \item \emph{Say it:} \emph{vā̃sḷī}
  \item \emph{Say it:} \emph{vā̃sḷī}, once more
  \item \emph{Say it:} \emph{mẽ kāle vā̃sḷī vagāḍī} — I played the flute yesterday.
  \item \emph{Recall:} say the Gujarati for song, then the Gujarati for to play an instrument, then say \emph{vā̃sḷī} again
\end{itemize}

\begin{tcolorbox}[breakable,colback=teal!4,colframe=teal!35!black,title={Before you move on}]
What does \gu{વાંસળી} mean? (\textquotedblleft{}Flute\textquotedblright{}.) Say it once more.
\end{tcolorbox}
\section[\texorpdfstring{\gu{તબલું} (tablũ)}{તબલું (tablũ)}]{\gu{તબલું} (tablũ) — tabla drum}

\label{lesson:GU-C254-tablu}



\begin{quote}
Before the new one: say the Gujarati for to play an instrument, then the Gujarati for flute.
\end{quote}

\subsection*{You'll want to know: \gu{તબલું}}
\textbf{\gu{તબલું}} — \emph{tablũ} — \textquotedblleft{}tabla drum\textquotedblright{}.

\begin{quote}
\textbf{\gu{મારા} \gu{ભાઈએ} \gu{તબલું} \gu{વગાડ્યું}.} — \emph{mārā bhāīe tablũ vagāḍyũ} — \textquotedblleft{}My brother played the tabla.\textquotedblright{}
\end{quote}

The pair is \textbf{\gu{તબલાં}}, \emph{tablā̃}.

\begin{cousinweb}
\textbf{\gu{તબલું}} (\emph{tablũ}) is a loan from Arabic \emph{ṭabl}, \textquotedblleft{}drum\textquotedblright{}, which reached Gujarati through Persian. Its Arabic root is \emph{ṭ-b-l}, \textquotedblleft{}to beat a drum\textquotedblright{} (Lane).
\end{cousinweb}

\begin{grammarlens}[title={What you've built}]
One more word for everyday life at A2.
\end{grammarlens}

\subsection*{Your turn}
\begin{itemize}
\raggedright
  \item \emph{Say it:} \emph{tablũ}
  \item \emph{Say it:} \emph{tablũ}, once more
  \item \emph{Say it:} \emph{mārā bhāīe tablũ vagāḍyũ} — My brother played the tabla.
  \item \emph{Recall:} say the Gujarati for to play an instrument, then the Gujarati for flute, then say \emph{tablũ} again
\end{itemize}

\begin{tcolorbox}[breakable,colback=teal!4,colframe=teal!35!black,title={Before you move on}]
What does \gu{તબલું} mean? (\textquotedblleft{}Tabla drum\textquotedblright{}.) Say it once more.
\end{tcolorbox}
\section[(dialogue)]{Words from chapters 249–254 — a review}

\label{lesson:GU-R254-review-249-254}



\begin{quote}
Say the Gujarati for mobile phone and number.
\end{quote}

\subsection*{Your turn}
Say these aloud:

\begin{itemize}
\raggedright
  \item mobile phone, number, message, voice, to obtain
  \item computer, screen, to press, to keep safe, to print
  \item meaning, translation, dictionary, to pay attention, to memorize
  \item duty, responsibility, principal, uniform, to retire
  \item hobby, leisure, to enjoy, to get bored, to try out
  \item music, song, to play an instrument, flute, tabla drum
\end{itemize}

\begin{tcolorbox}[breakable,colback=teal!4,colframe=teal!35!black,title={Before you move on}]
Mobile phone? (\textbf{\gu{મોબાઈલ}}.) Number? (\textbf{\gu{નંબર}}.) Message? (\textbf{\gu{સંદેશો}}.) Voice? (\textbf{\gu{અવાજ}}.) To obtain? (\textbf{\gu{મેળવવું}}.) Computer? (\textbf{\gu{કમ્પ્યૂટર}}.) Screen? (\textbf{\gu{સ્ક્રીન}}.) To press? (\textbf{\gu{દબાવવું}}.) To keep safe? (\textbf{\gu{સાચવવું}}.) To print? (\textbf{\gu{છાપવું}}.) Meaning? (\textbf{\gu{અર્થ}}.) Translation? (\textbf{\gu{અનુવાદ}}.) Dictionary? (\textbf{\gu{શબ્દકોશ}}.) To pay attention? (\textbf{\gu{ધ્યાન} \gu{આપવું}}.) To memorize? (\textbf{\gu{યાદ} \gu{કરવું}}.) Duty? (\textbf{\gu{ફરજ}}.) Responsibility? (\textbf{\gu{જવાબદારી}}.) Principal? (\textbf{\gu{આચાર્ય}}.) Uniform? (\textbf{\gu{ગણવેશ}}.) To retire? (\textbf{\gu{નિવૃત્ત} \gu{થવું}}.) Hobby? (\textbf{\gu{શોખ}}.) Leisure? (\textbf{\gu{નવરાશ}}.) To enjoy? (\textbf{\gu{માણવું}}.) To get bored? (\textbf{\gu{કંટાળવું}}.) To try out? (\textbf{\gu{અજમાવવું}}.) Music? (\textbf{\gu{સંગીત}}.) Song? (\textbf{\gu{ગીત}}.) To play an instrument? (\textbf{\gu{વગાડવું}}.) Flute? (\textbf{\gu{વાંસળી}}.) Tabla drum? (\textbf{\gu{તબલું}}.)
\end{tcolorbox}
